\section{Discussion}
\label{sec:discussion}

\subsection{Key Findings}

\textbf{Finding 1: Style-function dissociation as the primary mechanism.} Coverage without category analysis is misleading: pylint's 100\% total coverage of HumanEval failures overstates its diagnostic value. The 94.3\% C-category dominance means the LLM receives ``fix your formatting'' as its primary repair guidance on algorithmic failures. On complex tasks, this style-triggered rewriting corrupts logical structure, explaining the HumanEval regression. Future work reporting pylint coverage of LLM failures should decompose by flag category.

\textbf{Finding 2: Execution feedback is practically effective at fixed compute.} The +40.2pp MBPP improvement from a single round at $B$=1000 is substantively large---more than doubling baseline performance on a standard benchmark at minimal inference cost. The +4.9pp HumanEval improvement is smaller but robust. Execution feedback is a reliable choice for single-round repair across benchmark types.

\textbf{Finding 3: Task complexity moderates static analysis feedback utility.} Even low-information-content feedback (94.3\% style flags) helps on MBPP's simple tasks (+18.3pp). Style-guided rewrites preserve logical structure in short functions; they disrupt it in complex algorithms. The appropriate feedback signal may depend on the complexity profile of target tasks.

\subsection{Limitations}

\textbf{L1: Single model.} Results are based on Llama 3.1 8B Instruct only. Qwen2.5-Coder-7B replication was planned but not executed (resource constraints). Claims are scoped to Llama 3.1 8B.

\textbf{L2: Single repair round.} At $B$=1000, most problems complete one repair round. Results characterize single-round repair, not multi-round iterative improvement. Per-round data confirms this captures the primary effect.

\textbf{L3: P2 prediction refuted (informative null).} We predicted pylint coverage $<$50\%; actual coverage was 100\%. This null result on the primary metric is reported transparently. The functional coverage result (12.5\% E+W) is arguably more informative than the original prediction.

\textbf{L4: No per-problem qualitative analysis.} The ``style-guided corruption'' mechanism is inferred from aggregate results, not confirmed by per-problem solution comparison.

\subsection{Broader Impact}

This work contributes to responsible deployment of LLM-based code generation tools. Production systems using static analysis as repair feedback should be evaluated on functional correctness, not just coverage. The style-function dissociation may be invisible in aggregate quality metrics but visible in pass@k evaluations. No significant negative societal impacts are identified from this benchmarking study.
